
\documentclass{article}
\usepackage{colortbl}
\usepackage{makecell}
\usepackage{multirow}
\usepackage{supertabular}

\begin{document}

\newcounter{utterance}

\centering \large Interaction Transcript for game `matchit\_ascii', experiment `similar\_grid\_2', episode 0 with nemotron{-}3{-}ultra.
\vspace{24pt}

{ \footnotesize  \setcounter{utterance}{1}
\setlength{\tabcolsep}{0pt}
\begin{supertabular}{c@{$\;$}|p{.15\linewidth}@{}p{.15\linewidth}p{.15\linewidth}p{.15\linewidth}p{.15\linewidth}p{.15\linewidth}}
    \# & $\;$A & \multicolumn{4}{c}{Game Master} & $\;\:$B\\
    \hline

    \theutterance \stepcounter{utterance}  
    & & \multicolumn{4}{p{0.6\linewidth}}{
        \cellcolor[rgb]{0.9,0.9,0.9}{
            \makecell[{{p{\linewidth}}}]{
                \texttt{\tiny{[P1$\langle$GM]}}
                \texttt{Sodelujete v skupinski igri ugibanja.} \\
\texttt{Cilj je ugotoviti, ali sta ta mreža in druga slika, ki jo vidim samo jaz, enaki.} \\
\texttt{Najprej opišite svojo mrežo. Nato bom jaz podal svoj opis in si lahko drug drugemu zastavljava vprašanja o mrežah, da ugotoviva, ali sta enaki. Zdaj začnite svoj kratek opis mreže z "OPIS:", ki mu sledi opis. Ne dodajajte ničesar drugega. To je vaša mreža:} \\
\texttt{▢ X X X X} \\
\texttt{▢ X X X X} \\
\texttt{▢ X X X X} \\
\texttt{▢ X X X X} \\
\texttt{▢ ▢ ▢ ▢ ▢} \\
            }
        }
    }
    & \\ \\

    \theutterance \stepcounter{utterance}  
    & \multicolumn{4}{p{0.6\linewidth}}{
        \cellcolor[rgb]{0.9,0.9,0.9}{
            \makecell[{{p{\linewidth}}}]{
                \texttt{\tiny{[P1$\rangle$GM]}}
                \texttt{OPIS: Mreža je sestavljena iz 5 vrstic in 5 stolpcev. V prvih štirih vrsticah je prvi stolpec prazen (označen z ▢), ostalih štirje stolpci so polnjeni (označeni z X). Peta vrstica je v celoti prazna (vsi pet stolpcev so ▢).} \\
            }
        }
    }
    & & \\ \\

    \theutterance \stepcounter{utterance}  
    & & & \multicolumn{2}{p{0.3\linewidth}}{
        \cellcolor[rgb]{0.9,0.9,0.9}{
            \makecell[{{p{\linewidth}}}]{
                \texttt{\tiny{[GM$|$GM]}}
                \texttt{continue} \\
            }
        }
    }
    & & \\ \\

    \theutterance \stepcounter{utterance}  
    & & \multicolumn{4}{p{0.6\linewidth}}{
        \cellcolor[rgb]{0.9,0.9,0.9}{
            \makecell[{{p{\linewidth}}}]{
                \texttt{\tiny{[GM$\rangle$P2]}}
                \texttt{Sodelujete v skupinski igri ugibanja.} \\
\texttt{Cilj je ugotoviti, ali sta ta mreža in druga slika, ki jo vidim samo jaz, enaki.} \\
\texttt{Najprej opišite svojo mrežo. Nato bom jaz podal svoj opis in si lahko drug drugemu zastavljava vprašanja o mrežah, da ugotoviva, ali sta enaki. Zdaj začnite svoj kratek opis mreže z "OPIS:", ki mu sledi opis. Ne dodajajte ničesar drugega. To je vaša mreža:} \\
\texttt{X X X X X} \\
\texttt{▢ X X X ▢} \\
\texttt{▢ X X X X} \\
\texttt{▢ X X X X} \\
\texttt{▢ ▢ ▢ ▢ ▢} \\
            }
        }
    }
    & \\ \\

    \theutterance \stepcounter{utterance}  
    & & & \multicolumn{4}{p{0.6\linewidth}}{
        \cellcolor[rgb]{0.9,0.9,0.9}{
            \makecell[{{p{\linewidth}}}]{
                \texttt{\tiny{[GM$\langle$P2]}}
                \texttt{OPIS: 5x5 mreža. Prva vrstica: vsi X. Druga: prazno, tri X, prazno. Tretja in četrta: prazno, štiri X. Peta: vsi prazni.} \\
            }
        }
    }
     \\ \\

    \theutterance \stepcounter{utterance}  
    & & & \multicolumn{2}{p{0.3\linewidth}}{
        \cellcolor[rgb]{0.9,0.9,0.9}{
            \makecell[{{p{\linewidth}}}]{
                \texttt{\tiny{[GM$|$GM]}}
                \texttt{continue} \\
            }
        }
    }
    & & \\ \\

    \theutterance \stepcounter{utterance}  
    & & \multicolumn{4}{p{0.6\linewidth}}{
        \cellcolor[rgb]{0.9,0.9,0.9}{
            \makecell[{{p{\linewidth}}}]{
                \texttt{\tiny{[GM$\rangle$P2]}}
                \texttt{To je moj OPISOPIS: Mreža je sestavljena iz 5 vrstic in 5 stolpcev. V prvih štirih vrsticah je prvi stolpec prazen (označen z ▢), ostalih štirje stolpci so polnjeni (označeni z X). Peta vrstica je v celoti prazna (vsi pet stolpcev so ▢).} \\
\texttt{Zdaj postavi vprašanje, da bi ugotovil nove vidike moje mreže, ki se morda razlikujejo od tvoje mreže. Začni z "VPRAŠANJE:" in ne dodaj ničesar drugega.} \\
            }
        }
    }
    & \\ \\

    \theutterance \stepcounter{utterance}  
    & & & \multicolumn{4}{p{0.6\linewidth}}{
        \cellcolor[rgb]{0.9,0.9,0.9}{
            \makecell[{{p{\linewidth}}}]{
                \texttt{\tiny{[GM$\langle$P2]}}
                \texttt{VPRAŚVPRAŠĄȄ₁ₐ₁ₐₐₐ₁ₐₐₐₐₐ₁ₐₐ₁ₐₐ₁ₐₐ₁ₐₐ₁ₐₐ₁ₐₐ₁ₐₐ₁ₐₐ₁ₐₐ₁ₐₐ₁ₐₐ₁ₐₐ₁ₐₐ₁ₐₐ₁ₐₐ₁ₐₐ₁ₐₐ₁ₐₐ₁ₐₐ₁ₐₐ₁ₐₐ₁ₐₐ₁ₐₐ₁ₐₐ₁ₐₐ₁ₐₐ₁} \\
            }
        }
    }
     \\ \\

    \theutterance \stepcounter{utterance}  
    & & & \multicolumn{2}{p{0.3\linewidth}}{
        \cellcolor[rgb]{0.9,0.9,0.9}{
            \makecell[{{p{\linewidth}}}]{
                \texttt{\tiny{[GM$|$GM]}}
                \texttt{abort, first line: VPRAŚVPRAŠĄȄ₁ₐ₁ₐₐₐ₁ₐₐₐₐₐ₁ₐₐ₁ₐₐ₁ₐₐ₁ₐₐ₁ₐₐ₁ₐₐ₁ₐₐ₁ₐₐ₁ₐₐ₁ₐₐ₁ₐₐ₁ₐₐ₁ₐₐ₁ₐₐ₁ₐₐ₁ₐₐ₁ₐₐ₁ₐₐ₁ₐₐ₁ₐₐ₁ₐₐ₁ₐₐ₁ₐₐ₁ₐₐ₁ₐₐ₁ₐₐ₁, but expected VPRAŠANJE:} \\
            }
        }
    }
    & & \\ \\

\end{supertabular}
}

\end{document}
